% Generated draft: requires mathematical and source review.
\documentclass[11pt]{article}
\usepackage[margin=25mm]{geometry}
\usepackage{fontspec}
\setmainfont{Latin Modern Roman}
\usepackage{amsmath,amssymb,mathtools,unicode-math}
\setmathfont{Latin Modern Math}
\usepackage{enumitem,needspace,xurl,hyperref}
\setlength{\parindent}{0pt}
\setlength{\parskip}{5pt}
\setlength{\emergencystretch}{3em}
\newcommand{\N}{\mathbb{N}}
\newcommand{\Z}{\mathbb{Z}}
\newcommand{\Q}{\mathbb{Q}}
\newcommand{\R}{\mathbb{R}}
\newcommand{\C}{\mathbb{C}}
\newcommand{\F}{\mathbb{F}}
\newcommand{\A}{\mathbb{A}}
\newcommand{\PP}{\mathbb{P}}
\newcommand{\E}{\mathbb{E}}
\newcommand{\eps}{\varepsilon}
\newcommand{\dd}{\,\mathrm{d}}
\begin{document}
\section*{UnsolvedMath ID: 30001059}
\textbf{Source problem code:} catalogue entry
\subsection{Definitions and mathematical statement}
For each \(n\), let \(A_n=(c_{ij}\xi_{ij})_{1\leq i,j\leq n}\), where the coefficients are fixed with \(1\leq c_{ij}\leq C\), and \(C\) is bounded independently of \(n\). The \(\xi_{ij}\) are independent Bernoulli signs, conventionally taking \(1\) and \(-1\) with equal probability. Put \(\mu_n=\E|\det A_n|\).

The source describes a known estimate in the form
\[
|\det A_n|=\mu_n\exp(O(n^{2/3}))\quad\text{``almost surely,''}
\]
and asks whether the multiplicative loss can be reduced to \(n^{O(1)}\). The two-sided polynomial comparison event, under that reading, is
\[
\mu_n n^{-K}\leq |\det A_n|\leq\mu_n n^K,
\]
with \(K\) independent of \(n\) and of the coefficient array, allowing dependence on the fixed bound \(C\). The task is to obtain the intended probabilistic form of this event uniformly for the bounded heterogeneous arrays.

The probability qualifier is not precise in the extraction. Literal probability-one bounds for each finite \(n\), probability tending to one as \(n\to\infty\), and eventual almost-sure bounds for a coupled sequence are different assertions. No coupling across \(n\) or exceptional-event rate is supplied. The stated ``almost surely'' must therefore be checked rather than silently replaced by a high-probability claim. In particular a multiplicative lower comparison is sensitive to singular-matrix events. Likewise the source does not explicitly distinguish a two-sided comparison from a one-sided polynomial lower-loss statement.
\subsection{Short English statement}
For independent signed entries with magnitudes uniformly bounded above and below, seek polynomial rather than subexponential multiplicative control of the absolute determinant around its mean. The source's almost-sure qualifier and the exact tail formulation require clarification.
\subsection{Sources}
\begin{itemize}
\item Local UnsolvedMath registry, permanent ID \(30001059\), source code \texttt{catalogue entry}, ``Determinant Concentration for Heterogeneous Bernoulli Matrices.'' Registry curation: Ulam AI / UnsolvedMath Contributors, CC BY 4.0; underlying sources retain their supplied provenance.
\item \emph{Discrete Geometry} (2009), Oberwolfach Reports, pp.\ 2477--2556, DOI \url{https://doi.org/10.4171/owr/2008/44}.
\item Supplied ``Costello--Vu'': \url{https://arxiv.org/abs/0905.1909}.
\item Supplied literature triage, \texttt{literature\_checked\_at}: 2026-08-21: labeled open; says Costello--Vu improve concentration, including a stronger iid-case bound, but do not prove the conjectured polynomial loss for bounded heterogeneous Bernoulli entries. This is supplied dated curation, not independently verified; the exact probability estimates must be checked in the sources.
\end{itemize}
\end{document}
